\section{Projector geometry and finite Jacobi theory}
\label{sec:ensembles}

\subsection{Gauge-invariant response of a degenerate projector}

Let \(H(\lambda)\) have an isolated \(D\)-dimensional eigenspace at energy
\(E_0\), with spectral projector \(P\), complement \(Q=1-P\), and reduced
resolvent
\begin{equation}
 R=Q\bigl[Q(H-E_0)Q\bigr]^{-1}Q .
 \label{eq:reduced-resolvent}
\end{equation}
For the tangent \(V_\mu=\partial_\mu H\), define the
projector-to-complement response channel
\begin{equation}
 X_\mu=P V_\mu R .
 \label{eq:channel}
\end{equation}
In bases of the degenerate fiber and its coupled complement, \(X_\mu\) is a
\(D\times M\) matrix.
Kato perturbation theory gives, for two real directions \(v,w\),
\begin{equation}
 F_{vw}=i\left(X_vX_w^\dagger-X_wX_v^\dagger\right),\qquad
 \Gamma=X_vX_v^\dagger+X_wX_w^\dagger .
 \label{eq:curvature-metric}
\end{equation}
Here \(F_{vw}\) is Hermitian under our convention, and \(\Gamma\) is the
positive two-direction quantum metric.
Both transform by conjugation under a basis change within \(P\), so the
generalized eigenvalues of \((F_{vw},\Gamma)\) are gauge invariant.

The raw curvature scale depends on tangent normalization and energy
denominators.
We remove it by restricting to the \(r=\rank\Gamma\) support and defining
\begin{equation}
 \Omega=\Gamma^{-1/2}F_{vw}\Gamma^{-1/2}.
 \label{eq:normalized-curvature-intro}
\end{equation}
This normalization is algebraic and model agnostic.
It converts the diagnostic into the relative orientation of the two response
channels.
With
\begin{equation}
 B=\begin{pmatrix}X_v&X_w\end{pmatrix},\qquad
 J=\begin{pmatrix}0&iI_M\\-iI_M&0\end{pmatrix},
 \label{eq:BJ}
\end{equation}
whitening gives
\begin{equation}
 Y=\Gamma^{-1/2}B,\qquad
 YY^\dagger=I_r,\qquad
 \Omega=YJY^\dagger .
 \label{eq:isometry-compression}
\end{equation}
Thus \(\Omega\) is a compression of a fixed matrix with \(M\) positive and
\(M\) negative eigenvalues.
Its statistics depend only on an \(r\)-plane in
\(\mathbb C^{2M}\), i.e., a point on a complex Grassmannian.
This observation makes the random-matrix reference precise: randomness is
assigned exclusively to the orientation of the whitened response row space.

\subsection{Finite complex-Jacobi point process}

If the row space of \(Y\) is Haar distributed, the continuously distributed
eigenvalues of Eq.~\eqref{eq:isometry-compression} form a complex Jacobi
ensemble \cite{zyczkowski2000,collins2005}.
For \(r\leq M\), all \(k=r\) levels lie in \((-1,1)\) and
\begin{equation}
 p(\lambda_1,\ldots,\lambda_k)\propto
 \prod_{a<b}|\lambda_a-\lambda_b|^2
 \prod_{a=1}^{k}(1-\lambda_a^2)^{M-r}.
 \label{eq:jacobi-density}
\end{equation}
More generally, let
\begin{equation}
 k=\min(r,2M-r),\qquad a=|M-r| .
 \label{eq:jacobi-parameters}
\end{equation}
After the exact boundary levels derived below are removed, the interior joint
density has the same form as Eq.~\eqref{eq:jacobi-density} with \(r\) replaced
by \(k\) and exponent \(a\).

Let \(p_n^{(a)}(x)\) be orthonormal on \((-1,1)\) with weight
\(w_a(x)=(1-x^2)^a\), and define the weighted projection kernel
\begin{equation}
 {\cal K}_k(x,y)=\sqrt{w_a(x)w_a(y)}
 \sum_{n=0}^{k-1}p_n^{(a)}(x)p_n^{(a)}(y).
 \label{eq:jacobi-kernel}
\end{equation}
Its diagonal \(\rho_k(x)={\cal K}_k(x,x)\) integrates to \(k\).
The finite-rank unfolding is therefore
\begin{equation}
 u(x)=\int_{-1}^{x}\rho_k(s)\,\dd s .
 \label{eq:finite-unfolding}
\end{equation}
For the unfolded partition sum
\(Z_F(\tau)=\sum_{j=1}^{k}e^{-2\pi i\tau u(\lambda_j)}\), we use one
normalization throughout:
\begin{equation}
 \begin{aligned}
 K_{F,\mathrm{raw}}(\tau)
 &=\frac{1}{k}\left\langle |Z_F(\tau)|^2\right\rangle,\\
 K_{F,\mathrm{disc}}(\tau)
 &=\frac{1}{k}\left|\left\langle Z_F(\tau)\right\rangle\right|^2,\\
 K_{F,c}(\tau)&=K_{F,\mathrm{raw}}(\tau)
 -K_{F,\mathrm{disc}}(\tau).
 \end{aligned}
 \label{eq:sff-conventions}
\end{equation}
The determinantal identity
\(\rho_2(x,y)=\rho_k(x)\rho_k(y)-|{\cal K}_k(x,y)|^2\) then yields the exact
finite-\(k\) prediction
\begin{equation}
 K_{J,c}^{(k)}(\tau)
 =1-\frac{1}{k}\int_{-1}^{1}\!\!\dd x\int_{-1}^{1}\!\!\dd y\,
 e^{-2\pi i\tau[u(x)-u(y)]}|{\cal K}_k(x,y)|^2 .
 \label{eq:finite-jacobi-sff}
\end{equation}
Equation~\eqref{eq:finite-jacobi-sff} generates the finite-\(D\) ramp directly
and gives
\(K_{J,c}^{(k)}(0)=0\) from the projector identity and approaches the unit
plateau when the oscillatory cluster term is washed out.
We evaluate the double integral by Gauss--Jacobi quadrature; convergence and
an independent matrix-beta Monte Carlo comparison are given in
Appendix~\ref{app:numerics}.

The word ``form factor'' in Eq.~\eqref{eq:sff-conventions} denotes a Fourier
transform of the curvature point process, with \(\tau\) serving as a
dimensionless geometric Fourier coordinate.  A future dynamical extension
can connect this coordinate to operator growth and thermalization scales.

\subsection{Boundary-atom theorem and form-factor plateau}
\label{sec:atom-theorem}

When \(r>M\), the row space \(\cS\subset\mathbb C^{2M}\) must intersect both
\(M\)-dimensional eigenspaces \(\cH_\pm\) of \(J\):
\begin{equation}
 \dim(\cS\cap\cH_\pm)\geq r+M-2M=r-M .
 \label{eq:intersection-bound}
\end{equation}
Every vector in these intersections is an exact eigenvector of the
compression with eigenvalue \(\pm1\).
For a generic \(r\)-plane the bound is saturated, giving
\begin{equation}
 m_+=m_-=r-M,\qquad k=2M-r.
 \label{eq:atom-count}
\end{equation}
The atoms are deterministic under the Jacobi ensemble, so their connected
covariance is exactly zero.
If the full spectrum is normalized by its total rank \(r\), the exact
connected form factor is consequently
\begin{equation}
 K_{J,c}^{\mathrm{full}}(\tau)
 =\frac{k}{r}K_{J,c}^{(k)}(\tau),\qquad
 K_{J,c}^{\mathrm{full}}(\infty)=\frac{k}{r}.
 \label{eq:atom-plateau}
\end{equation}
This is an exact finite-dimensional plateau theorem.

For completeness, the raw, non-unfolded form factor retains deterministic
interference with the atoms.
Writing
\begin{equation}
 \begin{aligned}
 Z_{\rm full}(s)&=Z_{\rm at}(s)+Z_{\rm int}(s),\\
 Z_{\rm at}(s)&=m_-e^{2\pi i s}+m_+e^{-2\pi i s},
 \end{aligned}
 \label{eq:atom-partition}
\end{equation}
gives the exact decomposition
\begin{equation}
 K_{\rm raw}^{\rm full}
 =\frac{|Z_{\rm at}|^2}{r}
 +\frac{2\,\mathrm{Re}\!\left[
 Z_{\rm at}^{*}\langle Z_{\rm int}\rangle\right]}{r}
 +\frac{\langle|Z_{\rm int}|^2\rangle}{r}.
 \label{eq:atom-raw-decomposition}
\end{equation}
Equations~\eqref{eq:atom-plateau} and
\eqref{eq:atom-raw-decomposition} distinguish deterministic geometry from
correlated fluctuations and provide two independent numerical closure tests.
